\section{Introduction}
\label{sec:introduction}

\subsection{Motivation: a formal state object for persistent multi-agent cognition}
\label{subsec:intro-motivation}

Multi-agent systems that use language models, tools, and external data sources typically maintain memory, working context, and durable artifacts through stacks that combine retrieval, tool traces, and summarization \citep{yao2023react,qin2024toollearning}; long-horizon designs further separate working memory from external stores \citep{packer2024memgpt}.
Such stacks can be highly effective in deployment; they nonetheless leave open how to name a \emph{single} mathematical object that answers, at once: what thoughts exist, which are merely provisional, which are promoted to shared durable knowledge, and which transitions are permitted by the architecture itself.
This paper specifies such an object.

\emph{State of Thought} represents cognition as a \emph{layered directed typed multigraph}: vertices are thoughts; each carries a layer index, open or closed status, a persistence scope, and optional embedding; each edge has a type drawn from a finite contract (support, dependency, implication, contradiction, provenance, and other fixed labels).
The project label \emph{Matrioshka Brain} is informal; the mathematical object is \emph{State of Thought}.

\subsection{Why informal stacks are insufficient}
\label{subsec:intro-insufficient}

Scratch buffers, vector stores, and prompt templates can be combined ad hoc, but without an explicit state object, ``global memory'' and ``session context'' become implementation phrases rather than structural commitments.
Three distinctions must be primitive if persistence and multi-agent expansion are to be analyzed rather than merely implemented:
\begin{itemize}[leftmargin=*]
  \item \textbf{Cognitive state} at an index $t$: the full typed graph and labels (\Cref{sec:formal-model}).
  \item \textbf{Global state:} the subgraph of thoughts that are closed \emph{and} globally persistent---the durable cognitive core (\Cref{def:global-state}).
  \item \textbf{Session slice:} material retrieved from globals for a query, plus session-only open extensions (\Cref{def:session-slice}).
  \item \textbf{External interface:} observations admitted into layer~$0$ by ingest, before higher layers interpret them (\Cref{def:external-interface}).
\end{itemize}
Without this split, one cannot state precisely where exploration ends and commitment begins, nor separate evidence from abstraction in a uniform way.

\subsection{Design commitments: open and closed thought, layers, closure}
\label{subsec:intro-design}

An \emph{open} thought is provisional ($o(v)=\text{open}$); a \emph{closed} thought has been accepted for its layer and scope ($o(v)=\text{closed}$).
\emph{Closure} $\CloseOp_t^{\Pi}$ is the promotion gate: it applies admissibility, conflict policy, and persistence rules (\Cref{def:closure-op}).
Layers $\LayerSet=\{0,\ldots,H\}$ are not decorative: they index $\Vlayer{h}$ and enter admissibility (\Cref{def:admissibility}).
Layer~$0$ is observation and evidence; higher layers may link downward but do not rewrite layer~$0$ payload (\Cref{ax:layer-zero}).

\emph{Belief} (in the narrow sense used here) is a \emph{proper} subset of closed global thoughts: it additionally requires a layer threshold, admissibility, and non-blocked stance (\Cref{def:belief}).
Session-open hypotheses are not beliefs; layer~$0$ vertices are not beliefs.

\subsection{Contributions}
\label{subsec:intro-contributions}

This paper makes the following contributions.
\begin{enumerate}[leftmargin=*]
  \item \textbf{Object definition.} A complete specification of the State of Thought tuple, edge-type inventory, global and session substructures, and external ingest (\Cref{sec:formal-model}).
  \item \textbf{Operator signatures.} Session dynamics as maps---$\SliceOp$, $\StemOp$, $\ExpandOp$, $\MergeOp$, $\CloseOp$, $\IngestOp$---on graph-structured state (\Cref{sec:session-persistence}).
  \item \textbf{Policy-parameterized governance.} Conflict, merge, and closure decomposition through $\PsiConf$, $\PsiMerge$, and $\Pi=(\PsiConf,\SigmaAdm,\Pi_{\mathrm{prov}})$ (\Cref{def:conflict-policy,def:merge-policy,def:closure-op}), with conditional statements at stated epistemic status (\Cref{sec:theoretical}).
  \item \textbf{Reference instantiation.} A concrete triple $(\PsiConf^{\mathrm{ref}},\PsiMerge^{\mathrm{ref}},\Pi_{\mathrm{prov}}^{\mathrm{ref}})$ and a symbolic trace showing the definitions in action (\Cref{def:ref-conflict,def:ref-merge,def:ref-provenance}, \Cref{app:worked-example}).
\end{enumerate}

The central contribution is \emph{mathematical and architectural}: an explicit state object, policy-parameterized governance, and a clear split between definitions, the reference instantiation, and conjectures left for future analysis.

\subsection{Scope and epistemic status}
\label{subsec:intro-scope}

This is a pure theory-and-method manuscript: no software artifact, benchmark, or deployment metric is reported.
Literature positioning and scope relative to neighboring work appear in \Cref{sec:related}.
Where graph-theoretic consequences would require additional hypotheses, \Cref{sec:theoretical} and the appendix state them as conjectures or open questions rather than as finished theorems.

\subsection{Roadmap}
\label{subsec:intro-roadmap}

\Cref{sec:related} places the formalism next to blackboard-style shared memory, global-workspace analogies, recent agent-memory systems, and multilayer graph formalisms.
\Cref{sec:preliminaries} fixes typed multigraphs, layers, and status maps.
\Cref{sec:formal-model} defines the State of Thought object, admissibility, policies, closure, the reference triple, and an invariant classification table.
\Cref{sec:session-persistence} gives the end-to-end session pipeline and the reference merge policy.
\Cref{sec:theoretical} records definitional remarks, one conjecture, and pointers to further open questions.
\Cref{sec:discussion} unifies discussion with limitations and extensions; \Cref{sec:conclusion} summarizes and points forward.

\begin{figure}[t]
  \centering
  % Figure 1 — Global / Session / External overview (PROJECT_SPEC.md 11A)
% Layout pass: separate Session title from box; route arrows to avoid label collisions.
\resizebox{\linewidth}{!}{%
\begin{tikzpicture}[
  font=\small,
  nd/.style={circle, draw, minimum size=6mm, inner sep=0pt},
  open/.style={nd, dashed},
  arr/.style={-{Latex[length=2mm]}, thick},
]

\node[draw, rounded corners, minimum width=7.0cm, minimum height=3.55cm, align=left] (G) at (0,0) {};
\node[anchor=north west, inner sep=10pt] at (G.north west) {\textbf{Global State} (persistent closed thoughts)};
% Graph shifted down slightly for air under the heading
\foreach \i/\x/\y in {1/-1.65/-0.58,2/-0.55/-1.05,3/0.55/-0.42,4/1.55/-1.0,5/0.15/0.62} {
  \node[nd] (g\i) at ([xshift=\x cm, yshift=\y cm]G.center) {};
}
\draw[thick] (g1) -- (g2) -- (g3) -- (g5) -- (g1);
\draw[thick] (g2) -- (g4);

\node[draw=blue!55!black, rounded corners, minimum width=3.15cm, minimum height=1.95cm, thick] (S) at (5.05,0) {};
\node[anchor=south, font=\small\bfseries, yshift=3pt] at (S.north) {\textbf{Session State}};
\node[open] (o1) at ([xshift=-0.55cm,yshift=-0.18cm]S.center) {};
\node[open] (o2) at ([xshift=0.55cm,yshift=-0.38cm]S.center) {};
\node[nd] (c1) at ([xshift=0cm,yshift=0.52cm]S.center) {};
\draw[thick] (c1) -- (o1) -- (o2) -- (c1);

\node[draw, rounded corners, align=left, minimum width=2.25cm, inner sep=5pt] (E) at (-4.15,2.5) {\textbf{External}\\observations/tools};
% Enter global subgraph from the left at $g_1$ (avoids crossing the Global heading)
\draw[arr] (E.south) -- ++(0,-0.22) -| node[pos=0.18, below, sloped, font=\footnotesize, inner sep=1.5pt] {feeds} (g1.west);

\node[draw, rounded corners, inner sep=6pt] (Q) at (-4.05,-2.78) {\textbf{Query} $q$};
\draw[arr] (Q.east) -- ++(0.55,0) |- node[pos=0.58, below, font=\footnotesize, yshift=-2pt] {slice} (G.south west);

\draw[arr, dashed] (G.east) -- node[above, font=\footnotesize, yshift=3pt] {retrieve} (S.west);

\coordinate (clbend) at ($(S.east)+(1.2,0)$);
\draw[arr] (S.east) -- node[midway, above, font=\footnotesize] {closure} (clbend);
\draw[arr] (clbend) |- node[pos=0.38, right, font=\footnotesize, inner sep=2pt] {to global} (G.north east);

\node[anchor=west] at (-3.7,-3.2) {\emph{Open thoughts}: dashed nodes. \emph{Closed thoughts}: solid nodes.};

\end{tikzpicture}%
}

  \caption{Global persistent subgraph $\GlobalState_t$, session slice (globals plus session-only open vertices), external ingest to $\Vlayer{0}$, and closure write-back.
  Belief (\Cref{def:belief}) requires more than closure in the slice; stance $\stance_t$ omitted.}
  \label{fig:overview}
\end{figure}
